\newpage
\section{Wstęp}

\subsection{Tło problemu i cel pracy}

Praca dotyczy zaprojektowania i implementacji asystenta głosowego. Rozwiązanie to ma służyć osobom przebywającym w domkach wczasowych na terenie ośrodka wypoczynkowego Politechniki Warszawskiej w Wildze. Uczelnia prowadzi w domkach wczasowych badania, wprowadzając rozwiązania IoT, co prowadzi do powstania inteligentnej przestrzeni edukacyjnej typu \textit{living lab}. Głównym celem prac prowadzonych w ramach tej pracy magisterskiej będzie wdrożenie rozwiązania, które pozwoli wczasowiczom zdobyć szersze zrozumienie danych zbieranych w domku wczasowym. Interakcja z asystentem powinna być naturalna i płynna oraz intuicyjna dzięki jego głosowemu interfejsowi. Asystent zostanie oparty na otwartoźródłowych modelach LLM, a praca skupi się na porównaniu dostępnych technologii oraz ich analizie w kontekście wydajności pracy na sprzęcie z ograniczonymi zasobami obliczeniowymi.

Środowiskiem docelowym dla realizowanego projektu jest Ośrodek Wypoczynkowy Politechniki Warszawskiej w Wildze, modernizowany w ramach projektu Wilga MiLL. Miejsce to, dzięki wprowadzaniu rozwiązań Internetu Rzeczy (IoT) do domków wczasowych, łączy funkcję rekreacyjną z zapleczem technicznym typu \textit{living lab}. Wdrożona tam infrastruktura IoT funkcjonowała stabilnie przez długi czas, zapewniając ciągły monitoring parametrów środowiskowych oraz zużycia energii elektrycznej. Przez ten czas zespołowi badawczemu udało się zgromadzić wystarczającą ilość danych do analizy, które zostały wykorzystane w niniejszej pracy.

System czujników generuje znaczne ilości danych surowych, trudnych do analizy dla użytkownika systemu — wczasowicza. Choć są one bezcenne z perspektywy badawczej, dla przeciętnego wizytora, pozostają one niezrozumiałe i trudne do samodzielnej interpretacji. Przytłaczająca ilość informacji oraz brak intuicyjnego interfejsu mogą prowadzić do ignorowania potencjału systemu, a cel projektu, jakim jest zwiększenie świadomości o zużyciu energetycznym i wpływie człowieka na środowisko, staje się trudny do osiągnięcia bez odpowiedniego przetworzenia danych. Problem ten był już wcześniej obiektem badań zespołu badawczego. Aby zwiększyć dostępność informacji zbieranych przez system, zdecydowano się zaimplementować aplikację, która pomogłaby dostarczyć dane do użytkowników końcowych w sposób zrozumiały i intuicyjny. Niniejsza praca magisterska skupia się na rozwinięciu dostępnych interfejsów o asystenta głosowego, który podniósłby zainteresowanie systemem IoT. 

Istotnym aspektem warunków pracy projektowanego rozwiązania są również zachowania użytkowników. Dotychczasowe doświadczenia zebrane przez zespół badawczy w ośrodku wskazują na sceptyczny stosunek wczasowiczów do otaczającej ich w domkach technologii. Zauważono przypadki postrzegania urządzeń IoT jako zbędne, skomplikowane lub naruszające ich prywatność, co prowadziło do odłączania ich przez wczasowiczów od zasilania. Takie akcje podejmowane przez użytkowników w sposób oczywisty dezorganizują pracę systemu IoT. Wskazuje to na konieczność podniesienia atrakcyjności rozwiązania oraz wprowadzenia poczucia bycia aktywnym użytkownikiem systemu, a nie bycia obserwowanym przez czujniki. Interakcja z systemem nie może być dla wczasowicza obciążeniem, lecz atrakcyjnym elementem pobytu w Wildze. System powinien w sposób subtelny dostarczać ciekawych (dla osoby nietechnicznej) informacji, zachęcając do wejścia w interakcję z inteligentnym otoczeniem, zamiast prowokować do wyłączania węzłów systemu.

\subsection{Teza badawcza}

Z doświadczeń zebranych przez zespół badawczy wynika, że wczasowicze wykazują zainteresowanie systemem. Nie pozostają wobec niego obojętni, zwraca on ich uwagę. Problemem jest nieintuicyjność proponowanych rozwiązań. Kluczowe jest zbudowanie rozwiązania, dzięki któremu osoby będą chętnie edukowały się w zakresie ekologii oraz ich wpływu na środowisko. Dane zbierane przez system mają charakter techniczny, nie są one czytelne dla użytkownika końcowego. Zespół stale podejmuje wysiłki w celu podniesienia atrakcyjności \textit{living lab}. W pracy przyjęto założenie, że odpowiednie przedstawienie danych zbieranych przez system IoT może podnieść zainteresowanie proponowanym rozwiązaniem.


Na tej podstawie zbudowano tezę badawczą mówiącą, że asystent głosowy posiadający dynamicznie budowany kontekst oparty o dane czujnikowe może być atrakcyjnym elementem systemu. Co za tym idzie, może pomóc osobom nietechnicznym wejść w interakcję z systemem i zagłębić się w tematykę badań prowadzonych w Wildze. Interfejs głosowy powinien być naturalny oraz intuicyjnie przekazywać informacje, jego sposób odpowiedzi musi być przystępny dla osób nietechnicznych oraz sceptyków systemu.

Drugą płaszczyzną zawartą w tezie badawczej jest wymiar sprzętowy. Poczyniono założenie mówiące, że możliwe jest zbudowanie takiego rozwiązania w oparciu o platformę o bardzo ograniczonych zasobach. Założono pracę z Raspberry Pi 5. Ten mikrokomputer często uznawany jest w świecie IoT za sprzęt względnie szybki oraz mocny. Warto jednak pamiętać, że niniejsza praca magisterska dotyczy uruchamiania modeli LLM, zadanie to wymaga bardzo dużych zasobów obliczeniowych. W takich warunkach Raspberry Pi 5 wydaje się być platformą o bardzo mocno ograniczonych zasobach. Istotnym elementem pracy jest znalezienie balansu pomiędzy jakością odpowiedzi a czasem jej generowania.


\subsection{Zakres pracy}

Prace skupiają się zarówno na implementacji systemu asystenta głosowego jak i na przeprowadzeniu badań dotyczących uruchamiania modeli językowych na sprzęcie o ograniczonych zasobach. Celem implementacji jest zbudowanie systemu, który pozwoli dynamicznie budować kontekst dostarczany do asystenta oraz zbudowanie interfejsu głosowego asystenta, który umożliwi użytkownikom sprawną komunikację. Celem badań jest zrozumienie istotnych aspektów uruchamiania modeli LLM na Raspberry Pi 5, znalezienie optymalnego rozwiązania z punktu widzenia balansu między czasem odpowiedzi a jej jakością.

Prace podzielono na kilka etapów. Pierwszym z nich jest przegląd literatury. Ze względu na część badawczą pracy ten zakres prac uznano za istotny. Generatywna sztuczna inteligencja, której częścią są modele LLM, to bardzo dynamicznie rozwijająca się dziedzina. Zbadanie literatury w tym temacie pozwoli przeprowadzić badania w sposób zorganizowany oraz uporządkowany. 

Po dokonaniu analizy literatury przystąpiono do zbudowania koncepcji architektonicznej rozwiązania. Dzięki zbudowaniu stałych założeń węzłów systemu możliwe było jednóczesne rozwijanie systemu oraz prowadzenie badań nad jednym konkretnym węzłem- Raspberry Pi 5, na który optymalizowano prace modelu LLM. 

Następny etap pracy miał charakter badawczy. Skupiono się w nim na doborze i optymalizacji konfiguracji pracy modelu językowego uruchamianego na platformie o ograniczonych zasobach. Jako przykład takiej platformy stosowano oczywiście Raspberry Pi 5. Analizie poddano między innymi wpływ parametrów uruchomieniowych modelu, kwantyzacji i wielkości modelu, sposobu budowania promptu oraz warunków pracy urządzenia na jakość generowanej odpowiedzi. Prace skupiały się na ocenie  odpowiedzi w języku polskim, ponieważ z projektowanego rozwiązania będą korzystać docelowo osoby polskojęzyczne.


Po wykonaniu badań oraz opracowaniu koncepcji architektonicznych przystąpiono do implementacji systemu asystenta głosowego. System zbudowano z wykorzystaniem wniosków płynących z części badawczej oraz korzystając z założeń architektonicznych poczynionych we wcześniejszych etapach.  
Celem tego etapu było połączenie poszczególnych komponentów w spójny przepływ danych. Obejmował on przejście od głosowego zapytania użytkownika, przez przetworzenie wypowiedzi na tekst, zbudowanie promptu zapytania z dynamicznie budowanym kontekstem, uruchomienie lokalnego modelu językowego, aż po wygenerowanie odpowiedzi i jej odtworzenie w postaci mowy. W pracy uwzględniono również warstwę danych, ponieważ system IoT w Wildze nie był aktywny w trakcie prowadzenia prac. Zbudowano zatem warstwę symulującą przyspieszony upływ czasu, tak by kontekst zmieniał się szybciej niż w czasie rzeczywistym.

Ostatnim etapem było sformułowanie wniosków wynikających z przeprowadzonych testów. Część implementacyjna ograniczała się do budowania asystenta dla konkretnego celu- edukowania wczasowiczów w ośrodku w Wildze. Wnioski, które praca pozwoliła sformułować mają jednak charakter ogólny. Część badawcza skupiła się na problemie ogólnym, dotyczącym wielu zastosowań, co pozwoliło sformułować wnioski o charakterze ogólnym.
